\subsection{Caracterizaciones locales e infinitesimales de los morfismos lisos, de los morfismos no ramificados y de los morfismos étale.}
\label{subsection:IV.17.14}

\begin{proposition}[17.14.1]
\label{IV.17.14.1}
Sea $f:X\to Y$ un morfismo localmente de presentación finita, $x$ un punto de $X$, $y=f(x)$. Para que $f$ sea liso (resp. no ramificado, resp. étale) en el punto $x$, es necesario y suficiente que se satisfaga la condición siguiente:

Para todo esquema local $Y'=\operatorname{Spec}(A')$ con punto cerrado $y'$, todo morfismo $h:Y'\to Y$ tal que $h(y')=y$, todo subesquema cerrado $Y'_0=\operatorname{Spec}(A'/\mathfrak{J}')$ de $Y'$, donde el ideal $\mathfrak{J}'$ es de cuadrado nulo, y todo $Y$-morfismo $g_0:Y'_0\to X$ tal que $g_0(y')=x$, existe al menos un (resp. a lo sumo uno, resp. uno y solo uno) $Y$-morfismo $g:Y'\to X$ tal que $g_0$ sea su restricción a $Y'_0$.
\end{proposition}

Teniendo en cuenta las definiciones (17.1.1 y 17.3.1), se trata de demostrar que la condición del enunciado es \emph{suficiente} para que $f$ sea liso (resp. no ramificado) en el punto $x$.

\emph{(i) Caso de los morfismos lisos.} Se puede reducir al caso en que $Y=\operatorname{Spec}(A)$ y $X=\operatorname{Spec}(C)$ son afines, con $C=B/\mathfrak{R}$, donde $B=A[T_1,\ldots,T_n]$ es una $A$-álgebra de polinomios y $\mathfrak{R}$ un ideal de tipo finito de $B$. Para demostrar que $f$ es liso en el punto $x$, basta establecer que $\sh{O}_{X,x}$ es una $\sh{O}_{Y,y}$-álgebra formalmente lisa para las topologías discretas (17.5.1). Ahora bien, se tiene $\sh{O}_{X,x}=B_{\mathfrak q}/\mathfrak{R}B_{\mathfrak q}$, donde $\mathfrak q$ es un ideal primo de $B$, y, si $\mathfrak p$ es la imagen inversa de $\mathfrak q$ en $A$, $B_{\mathfrak q}$ es una $A_{\mathfrak p}$-álgebra formalmente lisa para las topologías discretas (0, 19.3.2 y 19.3.5). Por aplicación del criterio jacobiano (0, 22.6.1 y 20.5.12), basta, pues, ver que $B_{\mathfrak q}/\mathfrak{R}^2B_{\mathfrak q}$ es una extensión $A_{\mathfrak p}$-trivial de $B_{\mathfrak q}/\mathfrak{R}B_{\mathfrak q}$ por $\mathfrak{R}B_{\mathfrak q}/\mathfrak{R}^2B_{\mathfrak q}$. Pero esto resulta precisamente de la hipótesis aplicada a $A'=B_{\mathfrak q}/\mathfrak{R}^2B_{\mathfrak q}$ y $\mathfrak{J}'=\mathfrak{R}B_{\mathfrak q}/\mathfrak{R}^2B_{\mathfrak q}$ y al morfismo $g_0$ correspondiente al automorfismo identidad de $A'/\mathfrak{J}'=\sh{O}_{X,x}$.

\emph{(ii) Caso de los morfismos no ramificados.} Pongamos aquí $A=\sh{O}_{Y,y}$, $B=\sh{O}_{X,x}$, y observemos que, en virtud de (17.4.1, c')), basta demostrar que $\Omega^1_{B/A}=(\Omega^1_{X/Y})_x=0$. Ahora bien, con las notaciones de (0, 20.4.1), el anillo $C=\sh{P}^1_{B/A}=(B\otimes_A B)/\mathfrak{J}^2$ es local, puesto que lo es $C/(\mathfrak{J}/\mathfrak{J}^2)$, isomorfo a $B$. La hipótesis, aplicada a $A'=C$ y $\mathfrak{J}'=\mathfrak{J}/\mathfrak{J}^2$, muestra, por definición, que la aplicación $d_{B/A}:B\to\Omega^1_{B/A}$ es nula (0, 20.4.6), y, por tanto, que $\Omega^1_{B/A}=0$ por (0, 20.4.7).

\begin{proposition}[17.14.2]
\label{IV.17.14.2}
Sean $Y$ un preesquema localmente noetheriano, $f:X\to Y$ un morfismo localmente de tipo finito, $x$ un punto de $X$, $y=f(x)$. Para que $f$ sea liso (resp. no ramificado, resp. étale) en el punto $x$, es necesario y suficiente que satisfaga la condición de (17.14.1), suponiendo además que el anillo local $A'$ es artiniano, que $\mathfrak m'\mathfrak{J}'=0$, donde $\mathfrak m'$ es el ideal maximal de $A'$, y que el cuerpo residual $A'/\mathfrak m'$ de $A'$ es igual a $k(x)$.
\end{proposition}

Se trata nuevamente de demostrar que estas condiciones son suficientes en el caso de los morfismos lisos y en el de los morfismos no ramificados.

\oldpage[IV-4]{99}

\emph{(i) Caso de los morfismos lisos.} Si se pone $A=\sh{O}_{Y,y}$ y $B=\sh{O}_{X,x}$, se trata aquí de demostrar, teniendo en cuenta (17.5.3), que $B$ es una $A$-álgebra formalmente lisa para las topologías preádicas. Ahora bien, en virtud de (0, 22.1.4), esto se deduce de la hipótesis cuando se sustituye en ella la condición $\mathfrak m'\mathfrak{J}'=0$ por $\mathfrak{J}'\subset\mathfrak m'$. Pero, como $\mathfrak m'$ es nilpotente, se ve que la condición del enunciado ya es suficiente considerando los anillos $A'/\mathfrak m'^j\mathfrak{J}'$\footnote{El texto fuente escribe $\mathfrak m'^{,j}$ aquí y $\mathfrak m'^{,j}/\mathfrak m'^{,j+1}$ en (17.14.3); la coma espuria se elimina, como exigen la filtración por potencias y el testigo inglés mantenido.} y razonando por recurrencia sobre $j$, como en (0, 19.4.3).

\emph{(ii) Caso de los morfismos no ramificados.} Retomemos las notaciones de la demostración de (17.14.1, (ii)); para demostrar que el ideal $\mathfrak{J}/\mathfrak{J}^2$ de $C$ es nulo, observemos que $C$ es entonces un anillo local noetheriano, pues es un anillo local del anillo de un abierto afín de un subpreesquema de $X\times_YX$, que es localmente de tipo finito sobre $X$, y, por tanto, también sobre $Y$. Si $\mathfrak r$ es el ideal maximal de $C$, basta, pues, verificar que se tiene $\mathfrak{J}/\mathfrak{J}^2\subset\mathfrak r^n$ para todo $n$ (0$_{\mathrm I}$, 7.3.5), o también que $\mathfrak r^n=\mathfrak r^n+(\mathfrak{J}/\mathfrak{J}^2)$. Con las notaciones de (0, 20.4.6), esto significa también que, para todo $n$, las dos aplicaciones compuestas $B\xrightarrow{p_1}C\to C/\mathfrak r^n$ y $B\xrightarrow{p_2}C\to C/\mathfrak r^n$ son idénticas; pero, como $C/\mathfrak r^n$ es artiniano, esta identidad resulta de la hipótesis aplicada, por recurrencia sobre $n$, a $A'=C/\mathfrak r^n$ y $\mathfrak{J}'=(\mathfrak r^n+(\mathfrak{J}/\mathfrak{J}^2))/\mathfrak r^n$.

\begin{remark}[17.14.3]
\label{IV.17.14.3}
Obsérvese que, bajo las hipótesis de (17.14.2), si, además, el cuerpo $k(x)$ es una extensión \emph{finita} de $k(y)$, entonces los $A'$-módulos $\mathfrak m'^j/\mathfrak m'^{j+1}$ son $k(y)$-espacios vectoriales de dimensión finita, y, por tanto, \emph{a fortiori}, $A'$ es una $\sh{O}_{Y,y}$-álgebra finita.
\end{remark}
